\section{KinTCN}
\label{sec:method}

Computing a fitness-to-drive index requires a model that can extract
temporal patterns from multivariate driver time series in real time,
generalise across drivers with different baseline styles, and produce a
well-calibrated continuous score rather than a hard binary decision.
Temporal Convolutional Networks (TCNs)~\cite{Bai2018} satisfy all three
requirements: they enforce causality by construction, i.e., the output at time $t$
depends only on inputs at times $\leq t$, use dilated convolutions~\cite{vandenOord2016}
to capture dependencies spanning hundreds of timesteps without recurrence,
and are trained discriminatively end-to-end.

We propose \textbf{KinTCN}, a TCN trained on four standard in-vehicle kinematic
signals, i.e., speed ($v$), steering wheel angle ($\delta$), steering torque ($\tau$),
and lateral acceleration ($a_y$).
Internally, the network outputs a calibrated impairment probability
$\hat{p} \in [0,1]$; the fitness-to-drive index is its complement,
$f = 1 - \hat{p}$, so that $f$ close to 1 indicates a driver in full control
and $f$ close to 0 indicates elevated impairment risk.
The true network target, i.e., the driver's impairment state, is a latent variable that
cannot be observed directly. What can be observed are safety-critical error events, e.g., collisions, red-light
violations, recorded during the study.
KinTCN uses these events as \emph{weak supervision}: a kinematic window is
labelled as elevated-risk if an error is recorded in the near future, and as
safe otherwise. Such a supervision is weak as the label is a noisy proxy, i.e., an impaired driver may not
make an error in a given window, and an attentive driver may make one due to
road conditions, however it is the only supervision signal available without
intrusive physiological measurement.
The model therefore learns to score the driver's control state, not to predict
specific error events.
In our implementation, we paid careful attention to three sources of fragility that are common in
driver monitoring research: data leakage across participants during preprocessing,
class imbalance between rare error-preceding windows and abundant safe windows,
and overoptimistic evaluation on participants seen during training.
Each of these is addressed explicitly, and the design choices are detailed in
the subsections that follow.

\subsection{Label Derivation and Predictive Window}

The model operates on fixed-length kinematic windows and must assign each one
a scalar impairment probability before the next window arrives.
Because the driver's true impairment state is latent, each window is labelled
through the safety-critical events that follow it: a window ending at time $t$
is paired with the future interval
$[t + G,\; t + G + H]$,
and the events occurring within that interval constitute its weak supervision signal.

The gap $G$ enforces temporal separation between the kinematic observation
and the labelled period, preventing the model from exploiting signatures that
are contemporaneous with an error rather than preceding it.
The horizon $H$ defines the width of the future period over which evidence is
aggregated; a wider horizon increases label coverage at the cost of temporal
precision.
Both values are chosen empirically and their effect on performance is examined
in Section~\ref{sec:results}.

To reflect the differential safety relevance of each event, the dataset records
four types of safety-critical incident, each assigned a severity weight:
Collision ($w = 5$), Red-light violation ($w = 3$),
Panic braking with stop ($w = 2$), and Panic braking ($w = 1$).
The composite risk score for a window is then the sum of severity weights of
the distinct event types that occur at least once in the associated future interval.
A window is labelled positive ($y = 1$) if its score is $\geq 1$, i.e., if any
safety-critical event occurs in the future interval, and negative ($y = 0$) otherwise.
Higher-severity events increase the score without excluding lower-severity ones
from the positive class; the score is later used to up-weight high-severity
windows in the loss function (Section~\ref{sec:training}).

Since safety-critical incidents are rare relative to normal driving, the
labelling procedure inevitably produces a heavily skewed class distribution:
positive windows---those preceding an error---constitute a small minority of
the full window set.
This imbalance is structural rather than an artefact of the threshold choice:
even at the minimum score threshold of $1$, which maximises positive coverage by
counting any incident as sufficient, error-preceding windows remain far
outnumbered by safe ones.
A model trained on the raw distribution has little incentive to attend to the
minority class; it can minimise loss by assigning uniformly low impairment
scores, achieving high accuracy while being useless as a risk detector.
Correcting for this at both the data level and the loss level is therefore a
prerequisite for meaningful training, and the strategies adopted are described
in Section~\ref{sec:training}.

\subsection{Preprocessing and Feature Engineering}

Each driver exhibits a characteristic baseline style---habitual cruising speed,
typical steering amplitude, preferred torque range---that would dominate any
learned cross-participant representation if left uncontrolled.
Each kinematic signal is therefore z-scored per driving session (i.e.,
standardised to zero mean and unit variance) using the first $L + G$ seconds
of the session as the personal reference baseline, where $L$ denotes the
lookback duration (i.e., the length of the kinematic history fed to the
network at each prediction step).
This initial segment precedes the distraction induction and thus provides a
distraction-free snapshot of the driver's individual operating point.
Extreme values are clipped to a fixed number of standard deviations to suppress
artefacts from rare transient manoeuvres.
Because the model is evaluated under Leave-One-Participant-Out Cross-Validation (LOPO-CV),
these normalisation statistics must be recomputed in each fold using only the
training participants: allowing any statistic derived from the held-out driver
to influence preprocessing would introduce leakage and produce overoptimistic
generalisation estimates.

From each normalised recording, fixed-length windows of $L$ seconds are then
extracted with a uniform stride of 5\,s (i.e., consecutive windows overlap by
$L - 5$ seconds, so a new window starts every 5 seconds), converting the
continuous time series into a sequence of $(L \times 4)$ signal matrices.
Any window that extends beyond the end of the recording, or whose associated
future labelling interval $[t + G,\; t + G + H]$ is not fully contained within
the session, is discarded.
Notice that, the lookback duration $L$ controls how much temporal context the network can
attend to and, together with $G$ and $H$, is treated as a hyperparameter whose
effect on performance is examined in Section~\ref{sec:results}.

The raw four-channel window, however, provides the network with only
instantaneous signal values, giving it no direct access to how the signals are
evolving or how they deviate from local norms.
Each window is therefore transformed into a richer representation: for each of
the four kinematic signals, at every timestep, we compute the raw normalised
value, the first-order temporal difference, a within-window z-score, and
rolling mean and standard deviation at three temporal scales (short, medium,
long).
The multi-scale rolling statistics are motivated by the heterogeneous temporal
structure of the target events: the short scale responds to abrupt kinematic
changes characteristic of panic braking precursors, whereas the long scale
captures the gradual degradation of vehicle control that tends to precede
collisions and red-light violations.
The resulting feature tensor has shape $(L \times 36)$ per window
($4~\text{signals} \times 9~\text{features per signal}$).

\subsection{KinTCN: Model Design and Training}
\label{sec:training}

A central design requirement for KinTCN is that it must remain strictly causal:
at any moment $t$, the impairment score may only depend on signals observed up
to and including $t$, never on future measurements.
This rules out architectures that process the full sequence bidirectionally, such
as standard LSTMs or Transformers without causal masking.
KinTCN is instead built on Temporal Convolutional Networks
(TCNs)~\cite{Bai2018}, which enforce causality by construction through
one-sided, dilated convolutions that look only backwards in time.
Crucially, by stacking layers with exponentially increasing dilation (each
layer doubling the temporal spacing between the values it connects) the network
achieves a receptive field that covers the entire lookback window $L$ without
recurrence and without the vanishing-gradient problems that limit deep RNNs on
long sequences.

Not all moments within a window are equally informative: a brief episode of
erratic steering may be far more diagnostic than the surrounding normal driving,
yet a simple global average would dilute it.
After the convolutional stack, KinTCN therefore applies a lightweight temporal
attention mechanism that learns to weight each timestep according to its
relevance, producing a single summary vector that disproportionately reflects
the most diagnostically significant segments of the window.
This vector is then passed through a small fully connected network that produces
a single scalar: the raw impairment probability $\hat{p}$.

The class imbalance between error-preceding and safe windows requires
intervention before training can proceed meaningfully.
At the data level, we apply SMOTE~\cite{Chawla2002} (Synthetic Minority
Over-sampling Technique), which generates synthetic positive-class examples by
interpolating between existing real ones, moving the training distribution
closer to balance without discarding any safe-driving data.
To respect the sequential nature of the recordings, synthetic windows are only
constructed from real windows belonging to the same driving session and
sufficiently separated in time, ensuring that interpolated examples remain
temporally coherent.
At the loss level, we adopt focal loss~\cite{Lin2017}, which automatically
down-weights examples the model already handles confidently and concentrates the
training signal on ambiguous cases near the decision boundary, i.e., precisely
those windows where the kinematic evidence of impairment is subtle.
Each training sample is additionally weighted by its composite risk score, so
that windows preceding high-severity events such as collisions exert a stronger
pull on the model than those preceding minor incidents.

To reduce sensitivity to any single training run, $K$ independent models are
trained per fold with different random initialisations and their outputs averaged
into an ensemble score~\cite{Lakshminarayanan2017}.
Standard regularisation techniques (weight decay, additive noise on the input
signals, and random masking of short temporal segments) limit overfitting to
session-specific patterns, and training is halted early when performance on the
held-out validation fold ceases to improve.
Because raw neural network outputs are typically poorly calibrated, the ensemble
score is finally recalibrated using Platt scaling~\cite{Guo2017}, fitting a
logistic transform on the same validation fold.
The fitness-to-drive index delivered to the end user is the complement
$f = 1 - \hat{p}$, so that a score close to 1 reflects a driver in full
control and a score close to 0 signals elevated impairment risk.

\subsection{Evaluation Protocol}

\paragraph{Leave-One-Participant-Out Cross-Validation.}
All results are reported under Leave-One-Participant-Out Cross-Validation
(LOPO-CV): in each fold, one participant is held out entirely as the test set,
and all preprocessing statistics, scaler fitting, SMOTE generation, and model
training are performed exclusively on the remaining participants.
This is the strictest available cross-participant protocol and directly
simulates deployment on a driver never seen during development.

\paragraph{Evaluable and audit participants.}
Participants with fewer than five positive windows in their held-out fold
provide insufficient signal for reliable AUROC estimation and are designated
as \emph{audit drivers}.
For each audit driver, a dedicated LOPO fold is run---holding that driver
out entirely and training on all others---and the false positive rate (FPR)
at a fixed operating threshold is measured on their windows.
This separates performance estimation (evaluable drivers) from safety
auditing (audit drivers), providing a complete picture of model behaviour
across all participants.

\paragraph{Metrics.}
The primary metric is the area under the ROC curve (AUROC), which measures
the model's ability to rank impaired windows above safe ones independently
of any operating threshold.
Calibration is assessed by ECE~\cite{Guo2017} and Brier score.
Statistical significance relative to a gradient-boosted tree baseline
(XGBoost~\cite{Chen2016}) is assessed using the Wilcoxon signed-rank test
on per-participant AUROC values, with a two-sided alternative and
$\alpha = 0.05$.
Confidence intervals on aggregate metrics are obtained by
non-parametric bootstrap ($N = 2000$ resamples).
